\newpage
\section{Backend - Browsing}

\subsection{Panoramica}
\label{sec:browsing}

\noindent
\begin{tcolorbox}[colback=zpuslightgray,colframe=zpusgreen,title=Navigazione gerarchica,boxrule=1.5pt]
Il modulo Browsing implementa la navigazione gerarchica dei dati indicizzati in SQLite: 
DIP\ped{G} $\to$ Classi documentali $\to$ Processi $\to$ Documenti $\to$ File. 
Supporta inoltre getter alternativi per ID e stato di integrità.
\end{tcolorbox}

\subsubsection{Stack tecnologico}
\label{sec:browsing:stack}

\begin{table}[H]
  \centering
  \footnotesize
  \setlength{\tabcolsep}{4pt}
  \renewcommand{\arraystretch}{1.3}
  \begin{tabular}{|l|p{9cm}|}
    \hline
    \textbf{Componente} & \textbf{Tecnologia} \\
    \hline
    Runtime desktop & Electron (main + renderer) \\
    \hline
    Backend core & TypeScript \\
    \hline
    Dependency Injection & tsyringe \\
    \hline
    Database & better-sqlite3 \\
    \hline
    IPC frontend-backend & ipcMain.handle / ipcRenderer.invoke \\
    \hline
    Architettura & Clean Architecture (layering) \\
    \hline
  \end{tabular}
\end{table}

\subsubsection{Layer architetturali}
\label{sec:browsing:layers}

\begin{enumerate}
  \item \textbf{IPC Adapter:} boundary verso renderer, gestione canali, mapping DTO
  \item \textbf{Use Case:} thin orchestration layer, delega diretta a repository
  \item \textbf{Repository:} astrazione accesso dati, delega a DAO
  \item \textbf{DAO:} query SQL concrete, ricostruzione entity da righe
  \item \textbf{Mapper:} trasformazione entity $\leftrightarrow$ DTO
  \item \textbf{Entity / Value Objects:} modello di dominio
\end{enumerate}


\subsubsection{Organizzazione file backend}
\label{sec:browsing:struttura}

\begin{table}[H]
  \centering
  \footnotesize
  \setlength{\tabcolsep}{4pt}
  \renewcommand{\arraystretch}{1.5}
  \begin{tabular}{|l|p{10.5cm}|}
    \hline
    \textbf{Layer} & \textbf{File} \\
    \hline
    Entry point IPC & \texttt{main.ts}, \texttt{preload.ts}, \texttt{ipc-channels.ts}, \newline \texttt{BrowsingIpcAdapter.ts} \\
    \hline
    Use case & \texttt{dip/*}, \texttt{classe-documentale/*}, \newline \texttt{process/*}, \texttt{document/*}, \texttt{file/*} \\
    \hline
    Repository & \texttt{I*Repository.ts}, \texttt{*Repository.ts} \\
    \hline
    DAO & \texttt{I*DAO.ts}, \texttt{*DAO.ts}, \texttt{mappers/*Mapper.ts} \\
    \hline
    Entity / DTO & \texttt{entity/*}, \texttt{dto/*}, \texttt{value-objects/*} \\
    \hline
    Database & \texttt{schema.sql}, \texttt{DatabaseBootstrap.ts}, \texttt{IndexDip.ts} \\
    \hline
  \end{tabular}
\end{table}

\subsubsection{Flusso end-to-end}
\label{sec:browsing:flusso}


\paragraph{Sequenza runtime (passo a passo).}
\begin{enumerate}[labelwidth=0.8cm,leftmargin=1cm]
  \item \textbf{Renderer invoca IPC:} \texttt{window.electronAPI.invoke(channel, payload)}
  \item \textbf{Preload autorizza:} prefisso \texttt{ipc:*} consentito, inoltro a \texttt{ipcRenderer.invoke}
  \item \textbf{Main process intercetta:} \texttt{BrowsingIpcAdapter.register(ipcMain, dipPath)} registra handler
  \item \textbf{Risoluzione DI:} adapter usa container per risolvere il use case richiesto
  \item \textbf{Use case delega:} orchestra la chiamata al repository
  \item \textbf{Repository delega:} usa DAO per accedere ai dati
  \item \textbf{DAO esegue query:} SELECT su SQLite, ricostruzione entity da mapper
  \item \textbf{Conversione DTO:} adapter converte entity $\to$ DTO tramite \texttt{EntityToDtoConverter}
  \item \textbf{Risposta IPC:} renderer riceve DTO o null
\end{enumerate}

\paragraph{Flusso gerarchico di navigazione.}
\begin{center}
  \texttt{BROWSE\_GET\_DIP\_BY\_ID} \\
  $\downarrow$ \\
  \texttt{BROWSE\_GET\_DOCUMENT\_CLASS\_BY\_DIP\_ID} \\
  $\downarrow$ \\
  \texttt{BROWSE\_GET\_PROCESS\_BY\_DOCUMENT\_CLASS} \\
  $\downarrow$ \\
  \texttt{BROWSE\_GET\_DOCUMENTS\_BY\_PROCESS} \\
  $\downarrow$ \\
  \texttt{BROWSE\_GET\_FILE\_BY\_DOCUMENT}
\end{center}

Contratto IPC definito in \texttt{ipc-channels.ts} e registrato in \texttt{BrowsingIpcAdapter.ts}:

\begin{table}[H]
  \centering
  \footnotesize
  \setlength{\tabcolsep}{4pt}
  \renewcommand{\arraystretch}{1.4}
  \begin{tabular}{|p{4.2cm}|p{3.2cm}|p{2.8cm}|p{4.3cm}|}
    \hline
    \textbf{Canale} & \textbf{Payload} & \textbf{Return} & \textbf{Use case} \\
    \hline
    \texttt{BROWSE\_GET\_DOCUMENT\_BY\_ID}            & \texttt{id: number}                  & \texttt{DocumentDTO | null}      & \texttt{IGetDocumentByIdUC} \\
    \hline
    \texttt{BROWSE\_GET\_DOCUMENTS\_BY\_PROCESS}      & \texttt{processId: number}           & \texttt{DocumentDTO[]}           & \texttt{IGetDocumentByProcessUC} \\
    \hline
    \texttt{BROWSE\_GET\_DOCUMENTS\_BY\_STATUS}       & \texttt{status: IntegrityStatusEnum} & \texttt{DocumentDTO[]}           & \texttt{IGetDocumentByStatusUC} \\
    \hline
    \texttt{BROWSE\_GET\_FILE\_BY\_ID}                & \texttt{id: number}                  & \texttt{FileDTO | null}          & \texttt{IGetFileByIdUC} \\
    \hline
    \texttt{BROWSE\_GET\_FILE\_BUFFER\_BY\_ID}        & \texttt{id: number}                  & \texttt{Buffer | null}           & \texttt{IGetFileByIdUC} + lettura fs \\
    \hline
    \texttt{BROWSE\_GET\_FILE\_BY\_DOCUMENT}          & \texttt{documentId: number}          & \texttt{FileDTO[]}               & \texttt{IGetFileByDocumentUC} \\
    \hline
    \texttt{BROWSE\_GET\_FILE\_BY\_STATUS}            & \texttt{status: IntegrityStatusEnum} & \texttt{FileDTO[]}               & \texttt{IGetFileByStatusUC} \\
    \hline
    \texttt{BROWSE\_GET\_PROCESS\_BY\_ID}             & \texttt{id: number}                  & \texttt{ProcessDTO | null}       & \texttt{IGetProcessByIdUC} \\
    \hline
    \texttt{BROWSE\_GET\_PROCESS\_BY\_STATUS}         & \texttt{status: IntegrityStatusEnum} & \texttt{ProcessDTO[]}            & \texttt{IGetProcessByStatusUC} \\
    \hline
    \texttt{BROWSE\_GET\_PROCESS\_BY\_DOCUMENT\_CLASS} & \texttt{documentClassId: number}    & \texttt{ProcessDTO[]}            & \texttt{IGetProcessByDocumentClassUC} \\
    \hline
    \texttt{BROWSE\_GET\_DOCUMENT\_CLASS\_BY\_DIP\_ID} & \texttt{dipId: number}              & \texttt{DocumentClassDTO[]}      & \texttt{IGetDocumentClassByDipIdUC} \\
    \hline
    \texttt{BROWSE\_GET\_DOCUMENT\_CLASS\_BY\_STATUS} & \texttt{status: IntegrityStatusEnum} & \texttt{DocumentClassDTO[]}      & \texttt{IGetDocumentClassByStatusUC} \\
    \hline
    \texttt{BROWSE\_GET\_DOCUMENT\_CLASS\_BY\_ID}     & \texttt{id: number}                  & \texttt{DocumentClassDTO | null} & \texttt{IGetDocumentClassByIdUC} \\
    \hline
    \texttt{BROWSE\_GET\_DIP\_BY\_ID}                 & \texttt{id: number}                  & \texttt{DipDTO | null}           & \texttt{IGetDipByIdUC} \\
    \hline
    \texttt{BROWSE\_GET\_DIP\_BY\_STATUS}             & \texttt{status: IntegrityStatusEnum} & \texttt{DipDTO[]}                & \texttt{IGetDipByStatusUC} \\
    \hline
  \end{tabular}
  \caption{Contratto IPC del modulo Browsing}
  \label{tab:browsing:ipc}
\end{table}

\noindent
\textbf{Nota:} Il canale \texttt{BROWSE\_GET\_DIP\_BY\_DOCUMENT\_CLASS} è definito in \texttt{ipc-channels.ts} ma non registrato in \texttt{BrowsingIpcAdapter}.

\subsubsection{Use Case: firme getter}
\label{sec:browsing:usecase}

\noindent
\begin{tcolorbox}[colback=zpuslightgray,colframe=zpusgreen,title=Caratteristica della layer,boxrule=1pt]
Gli use case sono thin wrapper senza logica aggiuntiva; delegano direttamente al repository corrispondente.
\end{tcolorbox}

\paragraph{Interfacce per aggregato.}

\begin{table}[H]
  \centering
  \footnotesize
  \setlength{\tabcolsep}{4pt}
  \renewcommand{\arraystretch}{1.4}
  \begin{tabular}{|l|p{8.2cm}|}
    \hline
    \textbf{Use Case} & \textbf{Firma} \\
    \hline
    \texttt{IGetDipByIdUC} & \texttt{execute(id: number): Dip | null} \\
    \texttt{IGetDipByStatusUC} & \texttt{execute(status: IntegrityStatusEnum): Dip[]} \\
    \hline
    \texttt{IGetDocumentClassByDipIdUC} & \texttt{execute(dipId: number): DocumentClass[]} \\
    \texttt{IGetDocumentClassByIdUC} & \texttt{execute(id: number): DocumentClass | null} \\
    \texttt{IGetDocumentClassByStatusUC} & \texttt{execute(status: IntegrityStatusEnum): DocumentClass[]} \\
    \hline
    \texttt{IGetProcessByDocumentClassUC} & \texttt{execute(documentClassId: number): Process[]} \\
    \texttt{IGetProcessByIdUC} & \texttt{execute(id: number): Process | null} \\
    \texttt{IGetProcessByStatusUC} & \texttt{execute(status: IntegrityStatusEnum): Process[]} \\
    \hline
    \texttt{IGetDocumentByIdUC} & \texttt{execute(id: number): Document | null} \\
    \texttt{IGetDocumentByProcessUC} & \texttt{execute(processId: number): Document[]} \\
    \texttt{IGetDocumentByStatusUC} & \texttt{execute(status: IntegrityStatusEnum): Document[]} \\
    \hline
    \texttt{IGetFileByIdUC} & \texttt{execute(id: number): File | null} \\
    \texttt{IGetFileByDocumentUC} & \texttt{execute(documentId: number): File[]} \\
    \texttt{IGetFileByStatusUC} & \texttt{execute(status: IntegrityStatusEnum): File[]} \\
    \hline
  \end{tabular}
  \label{tab:browsing:usecase}
\end{table}

\subsubsection{Layer Repository}
\label{sec:browsing:repository}

\noindent
Ogni repository delega al DAO corrispondente il retrieval dei dati.

\begin{table}[H]
  \centering
  \footnotesize
  \setlength{\tabcolsep}{4pt}
  \renewcommand{\arraystretch}{1.5}
  \begin{tabular}{|l|l|p{7.2cm}|}
    \hline
    \textbf{Repository} & \textbf{Implementazione} & \textbf{Metodi getter} \\
    \hline
    \texttt{IDipRepository} & \texttt{DipRepository.ts} & 
      \texttt{ById}, \texttt{ByUuid}, \texttt{ByStatus} \\
    \hline
    \texttt{IDocumentClassRepository} & \texttt{DocumentClassRepository.ts} & 
      \texttt{ById}, \texttt{ByDipId}, \texttt{ByStatus} \\
    \hline
    \texttt{IProcessRepository} & \texttt{ProcessRepository.ts} & 
      \texttt{ById}, \texttt{ByDocumentClassId}, \texttt{ByStatus} \\
    \hline
    \texttt{IDocumentRepository} & \texttt{DocumentRepository.ts} & 
      \texttt{ById}, \texttt{ByProcessId}, \texttt{ByStatus} \\
    \hline
    \texttt{IFileRepository} & \texttt{FileRepository.ts} & 
      \texttt{ById}, \texttt{ByDocumentId}, \texttt{ByStatus} \\
    \hline
  \end{tabular}
  \label{tab:browsing:repositories}
\end{table}

\noindent
Inoltre: \texttt{save}, \texttt{updateIntegrityStatus} per persistenza; \texttt{searchDocument*}, 
\texttt{getIndexedDocumentsCount} per DocumentRepository.

\subsubsection{Layer DAO}
\label{sec:browsing:dao}

\noindent
Le interfacce \texttt{I*DAO} espongono metodi paralleli ai repository per il retrieval e query concrete eseguite in SQLite.

\begin{table}[H]
  \centering
  \footnotesize
  \setlength{\tabcolsep}{4pt}
  \renewcommand{\arraystretch}{1.5}
  \begin{tabular}{|l|p{2.5cm}|p{7.5cm}|}
    \hline
    \textbf{DAO} & \textbf{File} & \textbf{Query / Operazioni} \\
    \hline
    \texttt{DipDAO} & \texttt{DipDAO.ts} & 
      By id/uuid/status; UPSERT uuid \\
    \hline
    \texttt{DocumentClassDAO} & \texttt{DocumentClassDAO.ts} & 
      By id/dipId/status; search LIKE; UPSERT \\
    \hline
    \texttt{ProcessDAO} & \texttt{ProcessDAO.ts} & 
      By id/documentClassId/status; metadati albero \\
    \hline
    \texttt{DocumentDAO} & \texttt{DocumentDAO.ts} & 
      By id/processId/status; search avanzate; ricerca vettoriale \\
    \hline
    \texttt{FileDAO} & \texttt{FileDAO.ts} & 
      By id/documentId/status; sort is\_main DESC \\
    \hline
    \texttt{MetadataHelper} & \texttt{MetadataHelper.ts} & 
      Load/save ricorsivo; ricostruzione albero \\
    \hline
  \end{tabular}
  \label{tab:browsing:daos}
\end{table}

\subsubsection{Entity e Value Objects}
\label{sec:browsing:entity}

\noindent
Modello di dominio. Ogni entity espone getter e \texttt{setIntegrityStatus}:

\begin{table}[H]
  \centering
  \footnotesize
  \setlength{\tabcolsep}{4pt}
  \renewcommand{\arraystretch}{1.4}
  \begin{tabular}{|l|p{9cm}|}
    \hline
    \textbf{Entity} & \textbf{Attributi principali} \\
    \hline
    \texttt{Dip} & \texttt{id,uuid,integrityStatus} \\
    \hline
    \texttt{DocumentClass} & \texttt{id,dipId,dipUuid,uuid,name,timestamp,integrityStatus} \\
    \hline
    \texttt{Process} & \texttt{id,documentClassId,documentClassUuid,uuid,metadata,integrityStatus} \\
    \hline
    \texttt{Document} & \texttt{id,uuid,metadata,processId,processUuid,integrityStatus} \\
    \hline
    \texttt{File} & \texttt{id,uuid,filename,path,hash,isMain,documentId,documentUuid,integrityStatus} \\
    \hline
  \end{tabular}
  \label{tab:browsing:entities}
\end{table}

\subsubsection{Data Transfer Objects}
\label{sec:browsing:dto}

\noindent
Contratto dati IPC (restituiti dall'adapter al renderer):

\begin{table}[H]
  \centering
  \footnotesize
  \setlength{\tabcolsep}{4pt}
  \renewcommand{\arraystretch}{1.4}
  \begin{tabular}{|l|p{10cm}|}
    \hline
    \textbf{DTO} & \textbf{Attributi} \\
    \hline
    \texttt{DipDTO} & \texttt{id,uuid,integrityStatus} \\
    \hline
    \texttt{DocumentClassDTO} & \texttt{id,dipId,uuid,name,timestamp,integrityStatus} \\
    \hline
    \texttt{ProcessDTO} & \texttt{id,documentClassId,uuid,integrityStatus,metadata: MetadataDTO} \\
    \hline
    \texttt{DocumentDTO} & \texttt{id,processId,uuid,integrityStatus,metadata: MetadataDTO} \\
    \hline
    \texttt{FileDTO} & \texttt{id,documentId,filename,path,hash,integrityStatus,isMain} \\
    \hline
    \texttt{MetadataDTO} & \texttt{name,value (scalare|MetadataDTO[]),type: MetadataType} \\
    \hline
  \end{tabular}
  \caption{DTO schema (metadati caricati ricorsivamente)}
  \label{tab:browsing:dtos}
\end{table}

\noindent
\textbf{Nota:} Il converter genera eccezioni se \texttt{id} o FK sono \texttt{null} durante la conversione.


\subsubsection{Mapping e conversioni}
\label{sec:browsing:mapping}

\noindent
Trasformazioni tra layer di persistenza, entity e DTO:

\begin{table}[H]
  \centering
  \footnotesize
  \setlength{\tabcolsep}{4pt}
  \renewcommand{\arraystretch}{1.4}
  \begin{tabular}{|l|p{2.8cm}|p{8cm}|}
    \hline
    \textbf{Conversione} & \textbf{File} & \textbf{Funzioni} \\
    \hline
    Entity $\to$ DTO & \texttt{EntityToDtoConverter.ts} & 
      \texttt{*ToDto} per dip, docClass, process, document, file, metadata \\
    \hline
    SQL $\to$ Entity & \texttt{*Mapper.ts} & 
      \texttt{*Mapper} per dip, docClass, process, document, file, metadata \\
    \hline
  \end{tabular}
  \label{tab:browsing:mapping}
\end{table}

\noindent
\texttt{MetadataMapper.fromFlatRows} ricostruisce l'albero da righe flat con controllo cicli/root.

\subsubsection{Dependency Injection}
\label{sec:browsing:di}

\noindent
Configurazione in \texttt{container.ts}:

\begin{enumerate}
  \item \textbf{Provider DB:} \texttt{DATABASE\_PROVIDER\_TOKEN} $\to$ \texttt{DatabaseProvider}
  \item \textbf{Repository tokens:} \texttt{*\_REPOSITORY\_TOKEN} per dip, classe documentale, processo, documento, file
  \item \textbf{DAO tokens:} \texttt{*\_DAO\_TOKEN} per le implementazioni concrete
  \item \textbf{Use case tokens:} \texttt{GET\_BY\_*} per tutte le operazioni di retrieval
\end{enumerate}

\noindent
\texttt{BrowsingIpcAdapter} risolve ai root level: \texttt{DipUC}, \texttt{DocumentClassUC}, \texttt{ProcessUC}, \texttt{DocumentoUC}, \texttt{FileUC}.

\subsubsection{Schema database}
\label{sec:browsing:db}

\noindent
Tabelle coinvolte e indici in \texttt{schema.sql}:

\begin{table}[H]
  \centering
  \footnotesize
  \setlength{\tabcolsep}{4pt}
  \renewcommand{\arraystretch}{1.4}
  \begin{tabular}{|l|p{9cm}|}
    \hline
    \textbf{Tabella} & \textbf{Colonne chiave / Indici} \\
    \hline
    \texttt{dip} & PK \texttt{id}, AK \texttt{uuid}, IX \texttt{integrity\_status} \\
    \hline
    \texttt{document\_class} & PK \texttt{id}, FK \texttt{dip\_id}, AK \texttt{uuid}, \newline IX dip\_id + integrity\_status \\
    \hline
    \texttt{process} & PK \texttt{id}, FK \texttt{document\_class\_id}, AK \texttt{uuid}, \newline IX documentClass\_id + integrity\_status \\
    \hline
    \texttt{document} & PK \texttt{id}, FK \texttt{process\_id}, AK \texttt{uuid}, \newline IX process\_id + integrity\_status \\
    \hline
    \texttt{file} & PK \texttt{id}, FK \texttt{document\_id}, \newline UK (document\_id,path), IX ordinamento \\
    \hline
    \texttt{process\_metadata} & FK \texttt{process\_id}, albero \texttt{parent\_id} \\
    \hline
    \texttt{document\_metadata} & FK \texttt{document\_id}, albero \texttt{parent\_id} \\
    \hline
  \end{tabular}
  \label{tab:browsing:schema}
\end{table}

\subsubsection{Bootstrap e indicizzazione}
\label{sec:browsing:bootstrap}

\noindent
\textbf{DatabaseBootstrap.ts} (avvio):
\begin{enumerate}[labelwidth=0.8cm,leftmargin=1cm]
  \item Ricrea database zero
  \item Applica schema
  \item Chiama \texttt{IndexDip.execute(dipPath)}
\end{enumerate}

\noindent
\textbf{IndexDip.execute} (popolamento DB):
\begin{enumerate}[labelwidth=0.8cm,leftmargin=1cm]
  \item Insert: dip
  \item Insert: classi documentali
  \item Insert: processi (con metadati ricorsivi)
  \item Insert: documenti (con metadati ricorsivi)
  \item Insert: file
\end{enumerate}

\noindent
Browsing legge solo da queste tabelle; getter restituiscono \texttt{null}/\texttt{[]} se dati non presenti.


\subsubsection{Comportamenti runtime}
\label{sec:browsing:runtime}

\noindent
\begin{tcolorbox}[colback=zpuslightgray,colframe=zpusgreen,title=Semantica getter,boxrule=1pt]
\textbf{By-ID:} restituisce \texttt{null} se entity non trovata\\
\textbf{List:} restituisce array, anche vuoto\\
\textbf{DTO:} conversione applicata nel boundary IPC prima della restituzione
\end{tcolorbox}

\paragraph{Lettura file buffer (\texttt{BROWSE\_GET\_FILE\_BUFFER\_BY\_ID}).}
\begin{enumerate}[labelwidth=0.8cm,leftmargin=1cm]
  \item Recupera file entity tramite \texttt{IGetFileByIdUC}
  \item Risolve percorso assoluto: \texttt{path.resolve(dipPath, file.getPath())}
  \item Legge contenuto con \texttt{fs.readFileSync}
  \item In caso di errore I/O: log e ritorno \texttt{null}
\end{enumerate}

\paragraph{Sicurezza preload.}
\noindent
\texttt{preload.ts} consente IPC su prefissi autorizzati:\newline\texttt{ipc:search:}, \texttt{browse:}, \texttt{file:}, \texttt{check-integrity:}, \texttt{ipc:indexing:status}

\subsubsection{Test coverage}
\label{sec:browsing:test}

\noindent
\begin{table}[H]
  \centering
  \footnotesize
  \setlength{\tabcolsep}{4pt}
  \renewcommand{\arraystretch}{1.5}
  \begin{tabular}{|l|p{9cm}|}
    \hline
    \textbf{Suite} & \textbf{Copertura} \\
    \hline
    \texttt{BrowsingIpcAdapter.spec.ts} & 
      Canali; mapping DTO; semantica null \\
    \hline
    \texttt{*UseCases.spec.ts} & 
      Use case $\to$ repository \\
    \hline
    \texttt{*Repository.spec.ts} & 
      Repository $\to$ DAO \\
    \hline
    \texttt{*DAO.spec.ts} & 
      Query SQL; persistenza; ricostruzione entity \\
    \hline
  \end{tabular}
  \label{tab:browsing:tests}
\end{table}


\subsubsection{Osservazioni tecniche}
\label{sec:browsing:osservazioni}

\noindent
\begin{tcolorbox}[colback=zpuslightgray,colframe=zpusgreen,title=Issue identificati,boxrule=1pt]
\begin{enumerate}[labelwidth=0.8cm,leftmargin=1cm,nosep]
  \item \texttt{BROWSE\_GET\_DIP\_BY\_DOCUMENT\_CLASS} definito ma non registrato
  \item \texttt{ipc-channels.d.ts} non allineato a \texttt{.ts} (legacy costanti)
  \item Dichiarazioni duplicate in \texttt{IDocumentClassRepository}
  \item Registrazioni duplicate nel container DI
  \item Mapper attendono UUID FK: browsing OK, altri campi risky
  \item \texttt{electron-api.d.ts} vs \texttt{window.electronAPI.invoke} divergente
  \item DTO core vs renderer divergenti su metadati
\end{enumerate}
\end{tcolorbox}

\subsubsection{Raccomandazioni}
\label{sec:browsing:raccomandazioni}

\noindent
Backlog evolutivo:
\begin{enumerate}
  \item Allineare \texttt{ipc-channels.ts}, \texttt{.d.ts}, build artifacts
  \item Rimuovere duplicazioni interfacce + container DI
  \item Allineare query DAO ai campi richiesti dai mapper
  \item Consolidare DTO core $\leftrightarrow$ renderer
  \item Tipizzazione forte payload IPC per canale
\end{enumerate}

\newpage